% Simulazione d'esame: prova del 14/09/2026, esercizio 2 -- Mach in un ugello (schlieren).
% Stessi numeri e figure dello svolgimento di riferimento:
% svolgimenti/2026-09-14_E02_ugello_schlieren (calcoli.py e figure/).
% Compilare con: bash ../compila.sh 2026-09-14_E02_ugello_schlieren  (A4 + orologio)

% Preambolo comune delle simulazioni d'esame (vedi GUIDA.md).
% Lo stesso sorgente si compila in due formati:
%   - A4 (svolgimento.pdf), da leggere su carta, PC o telefono;
%   - orologio (orologio.pdf): colonna stretta, letta dall'app dell'orologio
%     (garmin-latex), che taglia le pagine in strisce larghe quanto lo schermo.
% Il formato orologio si sceglie definendo \formatoorologio prima di \input
% (lo fa compila.sh). In quel formato il margine destro è volutamente largo:
% ciò che sporge oltre la colonna (equazioni lunghe, tabelle, figure) viene
% riconosciuto dall'app e mostrato rimpicciolito, con lo zoom.

\documentclass[11pt]{article}

\usepackage[T1]{fontenc}
\usepackage[utf8]{inputenc}
\usepackage[italian]{babel}
\usepackage{amsmath,amssymb}
\usepackage{array}
\usepackage{siunitx}
\usepackage{booktabs}
\usepackage{graphicx}
\usepackage{float}
\usepackage{caption}
\usepackage{enumitem}
\usepackage{ragged2e}
\usepackage{titlesec}
\usepackage{hyperref}

\hypersetup{hidelinks}
\setlength{\parindent}{0pt}
\setlength{\parskip}{4pt}
\allowdisplaybreaks
\sisetup{output-decimal-marker={,}, group-separator={\,}, per-mode=symbol,
         exponent-product=\cdot, range-phrase={\text{--}}}

% Notazione del formulario (come template/svolgimento.tex).
\newcommand{\Cp}{C^{+}}
\newcommand{\Cm}{C^{-}}
\newcommand{\Mr}{M_{r}}
\newcommand{\Ms}{M_{s}}

% Separatore fra relazioni nella stessa formula: \begin{gather*} A \ee B \end{gather*}.
% In A4 le relazioni stanno sulla stessa riga, nel formato orologio una sotto l'altra.
\ifdefined\formatoorologio \newcommand{\ee}{,\\} \else \newcommand{\ee}{,\qquad} \fi

% Titolo dell'esercizio: \esercizio{Prova del GG/MM/AAAA -- Esercizio n}{argomento}.
% Va anche nei metadati del PDF: l'app dell'orologio lo usa come titolo nell'elenco.
\newcommand{\esercizio}[2]{%
  \title{#1\\[2pt]\large #2}\author{}\date{}%
  \hypersetup{pdftitle={#1: #2}}}

\ifdefined\formatoorologio
  % Colonna di 4,6 cm: a 2,4 px/pt è larga 262 px sullo schermo da 390.
  \usepackage[paperwidth=12cm,paperheight=150cm,left=0.2cm,top=0.2cm,
              textwidth=4.6cm,textheight=149cm,nohead,nofoot]{geometry}
  \usepackage[scaled=1.0]{helvet}
  \usepackage{sfmath}
  \renewcommand{\familydefault}{\sfdefault}
  \pagestyle{empty}
  \makeatletter\let\ps@plain\ps@empty\makeatother   % anche la pagina del titolo
  % colonna dichiarata per l'app dell'orologio (deve coincidere con left e textwidth)
  \hypersetup{pdfkeywords={colonna=0.2cm+4.6cm}}
  \RaggedRight
  \setlength{\parskip}{5pt}
  \hfuzz=20cm          % ciò che sporge va nel margine destro, senza avvisi
  \hbadness=10000
  \titleformat*{\section}{\normalsize\bfseries\RaggedRight\hyphenpenalty=10000}
  \titleformat*{\subsection}{\normalsize\bfseries\RaggedRight}
  \setlist{leftmargin=*,itemsep=2pt}
  \captionsetup{font=small,justification=raggedright,singlelinecheck=false}
  % figure: larghe il doppio della colonna, così sporgono e l'app le mostra con lo zoom
  \newcommand{\figura}[3][0.7]{%
    \begin{figure}[H]\includegraphics[width=9cm]{#2}\caption{#3}\end{figure}}
\else
  \usepackage[a4paper,margin=2.5cm]{geometry}
  \newcommand{\figura}[3][0.7]{%
    \begin{figure}[H]\centering\includegraphics[width=#1\textwidth]{#2}\caption{#3}\end{figure}}
\fi

\graphicspath{{../../svolgimenti/2026-09-14_E02_ugello_schlieren/figure/}}

\esercizio{Prova del 14/09/2026 -- Esercizio 2}{Mach nell'ugello a valle della gola}

\begin{document}
\maketitle

\emph{Traccia.} Stimare la distribuzione del numero di Mach nell'ugello a valle della gola. È data una
visualizzazione schlieren della regione di gola di un ugello (fig.~\ref{fig:ann}).

\section*{Che cosa mostra la figura}
La schlieren rende visibili i gradienti di densità: gli urti appaiono come linee nette e sottili, gli strati di
scorrimento come linee più sfumate e irregolari. Si vede, attraverso una finestra circolare, la regione di gola di un
ugello piano convergente--divergente: le pareti sono i blocchi neri in alto e in basso, la gola è la sezione più
stretta, in corrispondenza dei due fori, e la corrente va da sinistra verso destra. Le strutture sono nitide, quindi
tratto il campo come stazionario. L'uscita dell'ugello è fuori dalla figura e non ho né scale né pressioni: riferisco
le lunghezze all'altezza di gola $h^*$ e le pressioni alla pressione totale $p_0$, e assumo aria, $\gamma=1{,}4$.

\figura[0.62]{annotata.png}{Schlieren con le strutture che individuo. Ciano: gola; rosso: urti di separazione da W e
W$'$ e loro prosecuzione oltre l'incrocio P (tratteggio); arancio: strato separato inferiore; magenta: gambetto di
Mach T--T$'$ e urti forti. I numeri indicano gli stati.\label{fig:ann}}

La struttura che dice di più è la linea sottile che parte da ciascuna parete poco a valle della gola, in W e W$'$,
proprio dove lo strato limite si stacca e prosegue quasi parallelo all'asse. Una linea sottile che nasce dalla
parete è un urto obliquo, e da questo deduco tre cose:
\begin{enumerate}[label=(\alph*)]
\item la gola è efficace (strozzata);
\item a valle della gola la corrente è supersonica, perché solo lì può stare un urto obliquo;
\item la pressione all'uscita, che non vedo, è abbastanza alta da far separare la corrente dentro il condotto:
  l'ugello lavora sovraespanso.
\end{enumerate}
L'urto da W è della famiglia $\Cm$, quello da W$'$ della famiglia $\Cp$. Fra gli strati separati e le pareti, che
continuano ad aprirsi, restano le zone separate B e B$'$, ad aria quasi ferma. Le due linee si incrociano in P,
vicino all'asse, e proseguono fino al bordo destro della finestra, dove la luminosità cambia di colpo lungo un fronte
spezzato: ci torno alla fine. La tessitura a macchie e la fascia chiara lungo le pareti sono invece la grana della
stampa e un effetto ottico.

Il campo è quasi simmetrico, quindi mi basta una metà: per la separazione uso quella inferiore, dove lo strato
separato è netto, e per le strutture vicino all'asse la media delle due.

\section*{Modello}
Fino agli urti la corrente è adiabatica e isentropica; poiché è supersonica la studio con il modello 2D stazionario
supersonico (urti obliqui, onde di Mach, polari). Le zone separate le tratto come regioni a pressione uniforme
$p_B$. Il modello vale finché le regioni restano supersoniche: dietro gli urti forti alla fine della finestra non
vale più, e lì mi limito alle relazioni d'urto.

Prima di calcolare fisso il quadro del campo (fig.~\ref{fig:schema}) e gli stati che userò. Le pareti convesse
espandono la corrente a partire dalla linea sonica, più rapidamente vicino alle pareti: il nucleo 1 ha $M$ crescente
verso valle e verso le pareti. L'urto di separazione nasce forte a parete e si indebolisce verso l'asse; al termine
della zona visibile la compressione richiesta dalla contropressione avviene con un urto quasi normale (gambetto di
Mach), raccordato agli urti obliqui forti nei punti tripli, dietro cui la corrente è subsonica.

\figura[0.62]{schema.pdf}{Schema del campo in unità di $h^*$: pareti misurate (tratteggio: convergente), linea
sonica qualitativa, urti di separazione e loro prosecuzione oltre P, gambetto di Mach T--T$'$ e urti forti
(magenta), linee di scorrimento dai punti tripli (punteggiato), strati separati (arancio).\label{fig:schema}}

\begin{table}[H]
\centering\small
\begin{tabular}{cp{0.36\linewidth}p{0.36\linewidth}}
\toprule
Stato & Dove & Come lo si raggiunge \\
\midrule
1 & nucleo supersonico fra gola e urti (1$_W$ a parete prima di W$'$, sull'asse prima di T) & espansione isentropica \\
2, 2$'$ & dietro gli urti di separazione, lungo gli strati separati & urto di separazione \\
B, B$'$ & zone separate a parete, $p=p_B$ & strato separato \\
3 & dietro il gambetto di Mach, sull'asse & urto normale \\
4, 4$'$ & dietro gli urti forti da T, T$'$ & urto obliquo forte \\
\bottomrule
\end{tabular}
\caption{Stati del campo. Fra 3 e 4 (e fra 3 e 4$'$) c'è la linea di scorrimento che parte dal punto triplo.}
\end{table}

\section*{Dalla gola: il Mach non è uniforme sulla sezione}
Prima di misurare mi chiedo se basti un solo valore di $M$ per sezione. Vicino alla gola la teoria transonica di
Sauer--Meyer (piccole perturbazioni attorno a $M=1$, pareti con raggio di curvatura $R$) dà, con $\xi$ lungo l'asse,
$\eta$ trasversale ($\eta=\pm1$ a parete) e $d=1/\sqrt{(\gamma+1)R}$, la linea sonica e la sezione di gola
\begin{gather*}
\xi_{\rm sonica}=-\frac{(\gamma+1)d}{2}\,\eta^2 \ee \xi_{\rm gola}=-\frac d6\,\eta^2 .
\end{gather*}
A parete la linea sonica sta prima della gola, sull'asse dopo. Quindi, a parità di ascissa, la corrente vicino alla
parete ha già accelerato di più: $M$ è più alto a parete che sull'asse. Non serve la posizione esatta della linea
sonica; serve questa conseguenza: il Mach va stimato separatamente a parete e sull'asse.

\section*{A parete: il Mach alla separazione}
Misuro con il goniometro sulla metà inferiore, rispetto alla tangente alla parete nel piede W$'$
($x\approx0{,}21\,h^*$, parete inclinata di \ang{-15}):
\begin{itemize}
\item lo strato separato forma $\theta\approx\ang{14}$ con la tangente (è quasi parallelo all'asse);
\item l'urto di separazione forma $\sigma\approx\ang{59}$ con la stessa tangente (lo stesso valore lo leggo anche
  in W).
\end{itemize}
La tangente alla parete è la direzione della corrente subito a monte, quindi $\theta$ è la deviazione prodotta
dall'urto e $\sigma$ il suo angolo. Dalla relazione $\theta$--$\sigma$--$M$
\[
\tan\theta=2\cot\sigma\,\frac{M_1^2\sin^2\sigma-1}{M_1^2(\gamma+\cos2\sigma)+2}
\]
isolo $M_1^2$, moltiplicando per il denominatore e raccogliendo i termini in $M_1^2$:
\begin{gather*}
M_1^2=\frac{2(\tan\theta+\cot\sigma)}{2\cot\sigma\sin^2\sigma-\tan\theta\,(\gamma+\cos2\sigma)}
\ee M_1\approx1{,}62,
\end{gather*}
e dalla relazione isentropica $p_1/p_0=(1+\delta M_1^2)^{-\gamma/(\gamma-1)}=0{,}228$.

Controllo subito che l'urto sia sul ramo debole: per $M_1=1{,}62$ la deviazione massima è $\theta_{\max}=\ang{15.1}$,
raggiunta con $\sigma=\ang{65.7}$; ho $\theta<\theta_{\max}$ e $\sigma<\ang{65.7}$, quindi è il ramo debole, ma vicino
al distacco. Ho anche un controllo indipendente: lungo la parete, fino alla separazione, la corrente è un'onda
semplice partita dalla gola con $\theta=0$ e $\nu=0$, quindi $\nu(M)=|\theta_w|=\ang{14.9}$, che dà $M=1{,}60$.
I due valori coincidono entro la precisione del goniometro.

Dietro l'urto (stato 2$'$), con la componente normale $M_{n1}=M_1\sin\sigma=1{,}39$:
\begin{gather*}
\frac{p_2}{p_1}=1+\frac{2\gamma}{\gamma+1}\left(M_{n1}^2-1\right)=2{,}08 \ee
M_{n2}^2=\frac{1+\delta M_{n1}^2}{\gamma M_{n1}^2-\delta} \ee M_{n2}=0{,}745,\\
M_2=\frac{M_{n2}}{\sin(\sigma-\theta)}\approx1{,}05 \ee \frac{p_2}{p_0}\approx0{,}47 .
\end{gather*}
La corrente dietro l'urto è appena supersonica. Lo strato separato confina con la zona morta B, quindi
$p_B\simeq p_2\approx0{,}47\,p_0$. Il rapporto $p_1/p_B\approx0{,}48$ è dell'ordine di quelli tipici della
separazione negli ugelli sovraespansi (circa 0,4), quindi la lettura della figura è coerente.

\section*{Verso l'asse: un'onda di Mach}
Seguo l'urto di separazione verso l'asse. Vicino all'asse misuro, per la linea che parte da W$'$, \ang{44} prima
dell'incrocio P e \ang{43} dopo; per quella da W, \ang{41} e \ang{43}. In media circa \ang{43}, uguale prima e
dopo l'incrocio. Un urto di intensità finita, attraversando l'altro, entrerebbe
in una corrente già deviata e cambierebbe inclinazione; una linea che passa l'incrocio senza piegarsi è invece
un'onda di Mach. Allora l'angolo misurato è l'angolo di Mach del nucleo, e
\begin{gather*}
M=\frac{1}{\sin\mu}\approx1{,}47 \ee \frac{p}{p_0}=0{,}284 .
\end{gather*}
Questo è lo stato 1 sull'asse, da P fino alla fine della finestra ($x\approx0{,}8$--$1\,h^*$). È più basso del valore
a parete, anche se lo leggo più a valle: è proprio ciò che prevede Sauer--Meyer.

Quindi l'urto di separazione è rettilineo ma non ha intensità costante: a parete forma \ang{59} con la corrente,
sull'asse arriva come onda di Mach. La ragione è il nucleo non uniforme: verso l'asse $M$ è più basso e la corrente
è parallela all'asse. Per questo in P non c'è una vera interazione fra urti, e il nucleo sull'asse resta nello
stato 1.

\section*{In fondo alla finestra: gambetto di Mach e punto triplo}
\figura[0.24]{dettaglio_T.png}{Estremo destro della finestra: incrocio P, gambetto di Mach T--T$'$, urti forti.
\label{fig:det}}
Al bordo destro (fig.~\ref{fig:det}) le due linee, oltre P, arrivano su un tratto di urto perpendicolare all'asse, a
$x\approx0{,}99\,h^*$: è un \textbf{gambetto di Mach}, fra i \textbf{punti tripli} T e T$'$. Da lì il fronte piega
verso le pareti come \textbf{urto obliquo forte}, inclinato di circa \ang{79} sull'asse (\ang{81} da T, \ang{76} da
T$'$). Non lo considero il bordo della finestra: ha un tratto rettilineo spezzato che non segue il cerchio, e dietro
cambia la luminosità. È la ricompressione richiesta dalla contropressione.

Per semplicità assumo a monte lo stato 1 dell'asse: $M_1=1{,}47$, corrente parallela all'asse.

\emph{Gambetto (urto normale, 1$\to$3):}
\begin{gather*}
\frac{p_3}{p_1}=1+\frac{2\gamma}{\gamma+1}\left(M_1^2-1\right)=2{,}35 \ee M_3=0{,}71 \ee \frac{p_3}{p_0}=0{,}67 .
\end{gather*}

\emph{Urto forte con $\sigma=\ang{79}$ (1$\to$4$'$):} $M_{n1}=1{,}44$; dalla relazione $\theta$--$\sigma$--$M$ la
deviazione è $\theta\approx\ang{8}$, verso l'esterno. Poi
\begin{gather*}
\frac{p_4}{p_1}=2{,}26 \ee M_{n2}=0{,}722 \ee M_4=\frac{M_{n2}}{\sin(\sigma-\theta)}\approx0{,}76 \ee
\frac{p_4}{p_0}=0{,}64 .
\end{gather*}
Controllo il ramo: per $M_1=1{,}47$ si ha $\theta_{\max}=\ang{11.3}$ con $\sigma=\ang{66.9}$. Il mio urto ha
$\sigma=\ang{79}>\ang{66.9}$ e $\theta<\theta_{\max}$, quindi è sul ramo forte, e infatti $M_4<1$. Dietro tutto il fronte
la corrente è subsonica.

Il punto triplo è degenere: vi arriva un'onda di Mach, quindi gambetto e urto forte sono lo stesso fronte, quasi
normale vicino all'asse e più inclinato verso l'esterno. Infatti le pressioni dietro i due tratti ($0{,}67$ e
$0{,}64\,p_0$) sono quasi uguali. La piccola differenza la assorbe la curvatura del gambetto, che ho trascurato
assumendolo normale. La deviazione è la stima più fragile: per $M_1=1{,}47$ varia di circa mezzo grado per ogni grado
di $\sigma$, quindi $\ang{8}$ è indicativo.

\section*{La distribuzione del numero di Mach}
\figura[0.7]{mach.pdf}{Distribuzione di $M$ a valle della gola, a parete e sull'asse, con i valori dietro gli urti.
\label{fig:mach}}
Metto insieme i pezzi (fig.~\ref{fig:mach}):
\begin{itemize}
\item \textbf{a parete} $M$ vale 1 poco prima della gola e cresce rapidamente lungo la parete convessa (onda
  semplice, $\nu=|\theta_w|$) fino a circa 1,62 alla separazione; attraverso l'urto scende a circa 1,05;
\item \textbf{sull'asse} $M$ vale 1 poco dopo la gola e cresce più lentamente (tratto solo qualitativo) fino a
  circa 1,47 in P; resta tale fino al gambetto di Mach, dove scende a 0,71;
\item \textbf{dietro gli urti forti} $M\approx0{,}76$; nelle zone separate l'aria è quasi ferma.
\end{itemize}

\figura[0.62]{p_theta.pdf}{Piano $p$--$\theta$ della metà inferiore, pressione riferita a $p_0$.\label{fig:pth}}
Sul piano $p$--$\theta$ (fig.~\ref{fig:pth}) la curva punteggiata è la corrente attaccata alla parete, che arriva in
1$_W$ con $\theta=\ang{-15}$. La polare di 1$_W$ porta a 2$'$ ($\theta\approx\ang{-1}$) sul ramo debole, vicino al
massimo, e 2$'$ sta sulla retta $p=p_B$ delle zone separate. Sull'asse la polare dello stato 1 contiene sia il
gambetto (3, in cima, $\theta=0$) sia l'urto forte (4$'$, $\theta\approx\ang{-8}$), entrambi sul ramo forte e
molto sopra $p_B$. Le pressioni stanno fra $0{,}23$ e $0{,}67\,p_0$: basta la scala lineare.

\begin{table}[H]
\centering\small
\begin{tabular}{lcc}
\toprule
Posizione & $M$ & $p/p_0$ \\
\midrule
linea sonica & 1 & 0,528 \\
parete, alla separazione & $\approx1{,}62$ & 0,228 \\
asse, da P al gambetto & $\approx1{,}47$ & 0,284 \\
dietro l'urto di separazione & $\approx1{,}05$ & $\approx0{,}47$ \\
zone separate B, B$'$ & $\approx0$ & $\approx0{,}47$ \\
dietro il gambetto (3) & $\approx0{,}71$ & 0,67 \\
dietro gli urti forti (4, 4$'$) & $\approx0{,}76$ & 0,64 \\
\bottomrule
\end{tabular}
\caption{Riepilogo dei risultati.}
\end{table}

\section*{In conclusione}
Il modello ritrova tutto ciò che si vede: l'urto di separazione con l'angolo misurato, lo strato separato quasi
parallelo all'asse, le onde che attraversano l'asse senza piegarsi, il fronte quasi normale che si piega di circa
\ang{11} nei punti tripli. A valle di questo fronte la corrente è subsonica, a pressione più alta di $p_B$, e fuori
dalla figura riaccelera verso l'uscita. I limiti stanno nelle ipotesi: il nucleo non uniforme (qui solo stimato con
Sauer--Meyer), lo strato limite e la tridimensionalità della separazione, e la precisione delle misure d'angolo,
che conta soprattutto per l'urto forte.

\end{document}
